%BAB VIII: SOFTSKILL - KEMAMPUAN BEKERJASAMA DAN KOLABORASI
\section{MBKM -- Softskill: Kemampuan Bekerjasama dan Kolaborasi}

\begin{table}[H]
    \centering
    \begin{tabular}{ll}
        \textbf{Nama} & : Muhammad Argya Vityasy \\
        \textbf{NIM} & : 23/522547/PA/22475 \\
        \textbf{Pembimbing Akademik} & : Guntur Budi Herwanto, S.Kom., M.Cs. \\
        \textbf{Pembimbing Program MBKM} & : Guntur Budi Herwanto, S.Kom., M.Cs. \\
        \textbf{Mata Kuliah} & : Soft Skill: Kemampuan Bekerjasama dan Kolaborasi \\
        \textbf{Kode MK} & : MII21-4012 \\
        \textbf{Total SKS} & : 2 SKS \\
        \textbf{Judul Magang} & : MBKM Mandiri GDP Labs -- Digital Employee \\
        \textbf{Durasi} & : 9 Februari 2026 -- 22 Juni 2026 \\
    \end{tabular}
\end{table}

\subsection{Metodologi Kolaborasi}

Tim Software Engineering GDP Labs mengadopsi metodologi Agile/Scrum dalam pengembangan perangkat lunak. Setiap sprint biasanya berdurasi dua hingga empat minggu, dengan aktivitas rutin meliputi daily stand-up meeting (15 menit), sprint planning, sprint review, dan retrospective. Selain itu, setiap perubahan kode (fitur atau perbaikan) wajib dikirim melalui pull request (PR) di GitHub dan direview oleh setidaknya satu rekan tim sebelum dapat di-merge.

\subsection{Kegiatan}

\subsubsection{Daily Stand-up dan Sprint Planning}

Setiap pagi, tim engineering mengadakan daily stand-up untuk menyampaikan progress kemarin, rencana hari ini, serta hambatan yang dihadapi. Kegiatan ini memudahkan deteksi dini terhadap hambatan teknis (blockers) dan memungkinkan rekan tim untuk menawarkan bantuan secara proaktif. Selama sprint planning, dilakukan kolaborasi dengan Product Manager dan tim QA dalam memprioritaskan backlog dan menentukan estimasi kompleksitas tugas (story points).

\subsubsection{Code Review dan PR Workflow}

Setiap fitur yang dikembangkan dikirim melalui PR di GitHub dengan deskripsi yang jelas, mencakup konteks perubahan, daftar file yang dimodifikasi, dan screenshot atau log hasil pengujian jika relevan. Selain review manual oleh rekan tim, GDP Labs juga menggunakan AI-powered PR review agent (pr-agent) yang memberikan masukan awal terkait gaya kode, potensi bug, dan kesesuaian dengan konvensi proyek.

Dilakukan partisipasi aktif dalam memberikan dan menerima code review, yang membiasakan komunikasi teknis yang konstruktif, penerimaan terhadap kritik (growth mindset), serta pemahaman bahwa kode adalah tanggung jawab kolektif tim, bukan individu.

\subsubsection{GitHub Projects dan Task Tracking}

Proyek dikelola melalui GitHub Projects dengan kolom-kolom status: To Do, In Progress, In Review, Done. Selama magang, secara aktif diperbarui status tugas, menambahkan komentar pada issue terkait, serta menjelaskan alasan teknis setiap keputusan dalam thread diskusi. Fungsi GitHub Issues juga digunakan untuk melacak bug, perencanaan fitur, dan diskusi desain bersama tim.

\subsection{Refleksi}

Kolaborasi dalam lingkungan tim Agile memerlukan lebih dari sekadar kemampuan teknis; diperlukan komunikasi transparan, empati terhadap perspektif orang lain, dan kesadaran bahwa hambatan tim adalah prioritas bersama. Code review, meskipun sering kali menantang secara emosional, ternyata menjadi instrumen belajar yang sangat efektif untuk memahami kualitas kode dan best practices yang diterapkan rekan-rekan senior. Pengalaman ini memperkuat keyakinan bahwa perangkat lunak berkualitas tinggi dibangun oleh tim, bukan oleh individu tunggal.

\clearpage

